\documentclass[a4paper,12pt]{article}
\usepackage[T2A]{fontenc}
\usepackage[utf8]{inputenc}
\usepackage[russian]{babel}
\usepackage{amsmath,amssymb,mathtools}
\usepackage{helvet}
\renewcommand{\familydefault}{\sfdefault}
\usepackage{icomma}
\usepackage[a4paper,left=21mm,right=21mm,top=16mm,bottom=16mm,headheight=16pt]{geometry}
\usepackage{microtype}
\usepackage{tikz}
\usetikzlibrary{calc,arrows.meta,positioning,shapes.geometric}
\usepackage{tabularx,booktabs,array}
\usepackage{enumitem}
\usepackage[hidelinks]{hyperref}
\usepackage{parskip}
\usepackage{fancyhdr}
\usepackage{setspace}
\usepackage{xcolor}
\definecolor{accent}{HTML}{2E6F73}
\definecolor{darkaccent}{HTML}{18484B}
\definecolor{textgray}{HTML}{2E3133}
\definecolor{rulegray}{HTML}{B9CBCA}
\color{textgray}
\onehalfspacing
\setlength{\parindent}{0pt}
\setlength{\parskip}{0.18em}
\setlist[itemize]{leftmargin=1.5em,itemsep=0em,topsep=0.08em}
\setlist[enumerate]{leftmargin=1.6em,itemsep=0.30em,topsep=0.12em}
\pagestyle{fancy}
\fancyhf{}
\fancyhead[L]{\small\color{darkaccent}Равномерное движение по окружности}
\fancyhead[R]{\small\color{darkaccent}10.10.26}
\fancyfoot[C]{\small\thepage}
\renewcommand{\headrulewidth}{0.4pt}
\newcommand{\topic}[1]{\vspace{0.3em}{\large\bfseries\color{darkaccent}#1}\par\vspace{0.1em}{\color{rulegray}\rule{\linewidth}{0.5pt}}\vspace{0.05em}}
\newcommand{\term}[1]{\textbf{\color{darkaccent}#1}}
\newcommand{\eqbox}[1]{\[\boxed{#1}\]}
\newcommand{\unit}[1]{\,\text{#1}}

\begin{document}
\thispagestyle{empty}
\begin{center}
\vspace*{25mm}
{\Large\bfseries\color{accent}ЧЕК-ЛИСТ}\\[10mm]
{\Huge\bfseries\color{darkaccent}Равномерное движение\\[2mm] по окружности}\\[10mm]
{\large Формулы, сравнение вращений и центростремительное ускорение}\\[15mm]
{\color{accent}\rule{0.64\linewidth}{0.8pt}}\\[15mm]
{\Large Володя Хачатурян}\\[4mm]
{\large 10 октября 2026 года}\\[19mm]
\begin{minipage}{0.84\linewidth}\centering
Понять, как связаны период, частота, линейная и угловая скорости; научиться отличать движение точек одного диска от вращения сцепленных шестерён.
\end{minipage}
\end{center}
\vfill
\begin{center}
{\small\color{darkaccent}Лёвин Артём Александрович эксклюзивно для Володи Хачатуряна}
\end{center}
\vspace{6mm}
\begin{tikzpicture}[remember picture,overlay]
\draw[accent,line width=1.1pt] ($(current page.south west)+(18mm,7mm)$)
 .. controls ($(current page.south west)+(59mm,12mm)$) and ($(current page.south east)+(-50mm,3mm)$)
 .. ($(current page.south east)+(-18mm,7mm)$);
\end{tikzpicture}
\clearpage

\topic{1. Период и частота}
Пусть точка совершила $n$ оборотов за время $t$. Число оборотов \emph{может быть дробным}: четверть окружности соответствует $n=\frac14$.
\eqbox{\nu=\frac{n}{t},\qquad T=\frac{t}{n},\qquad \nu=\frac1T}
\begin{itemize}
 \item $T$ --- \term{период}: время одного оборота (с).
 \item $\nu$ --- \term{частота}: число оборотов за секунду (Гц = с$^{-1}$).
 \item Из формул: $n=\nu t=\dfrac{t}{T}$; $t=nT=\dfrac n\nu$.
\end{itemize}
\textbf{Пример.} За $1\unit{с}$ точка прошла четверть окружности. Тогда
\[\nu=\frac{1/4}{1}=0{,}25\unit{Гц},\qquad T=\frac1{0{,}25}=4\unit{с}.\]
\textbf{Смысл:} один полный оборот при неизменной скорости занял бы $4\unit{с}$.

\topic{2. Линейная и угловая скорости}
При равномерном движении по окружности модуль линейной скорости $v$ постоянен, хотя направление вектора скорости меняется. За оборот пройден путь $2\pi R$.
\eqbox{v=\frac{2\pi Rn}{t}=2\pi R\nu=\frac{2\pi R}{T}}
Угловая скорость $\omega$ показывает, насколько быстро изменяется угол поворота $\varphi$, измеренный \term{в радианах}. Один оборот: $2\pi\unit{рад}=360^\circ$.
\eqbox{\omega=\frac{\varphi}{t}=\frac{2\pi n}{t}=2\pi\nu=\frac{2\pi}{T}}
\eqbox{v=\omega R}
\textbf{Размерности:} $R$ --- м; $v$ --- м/с; $\omega$ --- рад/с. В формулах с углом используйте радианы.

\topic{3. Центростремительное ускорение}
Даже при постоянном модуле скорости направление движения меняется. Вектор $\vec a_{\mathrm c}$ в каждой точке направлен \term{к центру окружности}.
\eqbox{a_{\mathrm c}=\frac{v^2}{R}=\omega^2R}
\textbf{Быстрый анализ формулы:} при неизменном $R$ увеличение $v$ в $3$ раза увеличивает $a_{\mathrm c}$ в $9$ раз. При неизменном $v$ увеличение $R$ в $5$ раз уменьшает $a_{\mathrm c}$ в $5$ раз.

\clearpage
\topic{4. Две принципиально разные ситуации}
\textbf{Ситуация А. Две точки одного твёрдого вращающегося диска.} За одинаковое время они поворачиваются на одинаковый угол. Поэтому
\eqbox{\omega_A=\omega_B,\quad \nu_A=\nu_B,\quad T_A=T_B,\qquad \frac{v_A}{v_B}=\frac{R_A}{R_B}}
Точка дальше от оси проходит больший путь за то же время, а значит, имеет больший модуль линейной скорости.

\begin{center}
\begin{tikzpicture}[scale=0.92,>=Latex]
\draw[accent,thick] (0,0) circle (1.7);
\fill[darkaccent] (0,0) circle (1.6pt) node[below left] {$O$};
\draw[accent!80,thick] (0,0)--(1.50,0);
\fill[darkaccent] (0.85,0) circle (2pt) node[above] {$B$};
\fill[darkaccent] (1.50,0) circle (2pt) node[above] {$A$};
\draw[->,accent,thick] (1.50,0)--(1.50,0.8) node[right] {$\vec v_A$};
\draw[->,accent,thick] (0.85,0)--(0.85,0.45) node[left] {$\vec v_B$};
\node[align=left,anchor=west,text width=7.9cm] at (2.20,0.15) {Один диск: $\omega_A=\omega_B$, но при $R_A>R_B$ имеем $v_A>v_B$.};
\end{tikzpicture}
\end{center}
\textbf{Пример.} На одном диске $R_A=20\unit{см}$, $R_B=15\unit{см}$; $v_B=1{,}5\unit{м/с}$. Тогда
\[v_A=v_B\frac{R_A}{R_B}=1{,}5\cdot\frac{20}{15}=2\unit{м/с}.\]

\textbf{Ситуация Б. Две шестерни с внешним зацеплением без проскальзывания.} Модули \term{линейных скоростей зубьев в месте контакта} равны, а шестерни вращаются в противоположных направлениях.
\eqbox{v_1=v_2,\quad R_1\omega_1=R_2\omega_2,\quad \frac{\nu_1}{\nu_2}=\frac{R_2}{R_1}}
Радиус больше $\Rightarrow$ частота меньше. \textbf{Не путайте:} у двух точек одного диска одинаковы $\omega$, $\nu$, $T$; у двух зацепленных шестерён равны модули линейных скоростей \emph{на ободах}.

\textbf{Пример.} Первая шестерня радиусом $20\unit{см}$ сделала $20$ оборотов за $10\unit{с}$. Вторая имеет радиус $10\unit{см}$. Тогда
\[\nu_1=\frac{20}{10}=2\unit{Гц},\qquad \nu_2=\nu_1\frac{R_1}{R_2}=2\cdot\frac{20}{10}=4\unit{Гц}.\]

\topic{5. Практическая памятка и работа с формулами}
\begin{itemize}
 \item $R=d/2$. Всегда согласуйте единицы: $1\unit{см}=10^{-2}\unit{м}$, $1\unit{км}=10^3\unit{м}$.
 \item В школьных приближениях: сутки $=24\cdot3600=86\,400\unit{с}$; год $=365\cdot24\cdot3600=31\,536\,000\unit{с}$.
 \item При делении степеней десятки показатели вычитаются: $10^9/10^3=10^6$.
 \item При поиске неизвестной величины сначала \term{выразите её из формулы}, затем подставьте числа: $T=t/n\Rightarrow n=t/T$; $\nu=n/t\Rightarrow t=n/\nu$.
 \item \textbf{Оценка для орбиты Земли:} при $R\approx1{,}5\cdot10^{11}\unit{м}$ и $T\approx3{,}1536\cdot10^7\unit{с}$ скорость $v=2\pi R/T\approx30\unit{км/с}$.
\end{itemize}

\clearpage
\begin{center}{\LARGE\bfseries\color{darkaccent}Алгоритм: что одинаково при вращении?}\end{center}
\vspace{8mm}
\begin{center}
\begin{tikzpicture}[
 >=Latex,
 start/.style={ellipse,draw=accent,thick,text width=6.0cm,align=center,minimum height=12mm},
 choice/.style={diamond,draw=accent,thick,aspect=2.4,text width=3.6cm,align=center,inner sep=1.0pt},
 action/.style={rounded corners=2pt,draw=darkaccent,thick,text width=5.2cm,align=center,minimum height=15mm,inner sep=2mm},
 note/.style={rounded corners=2pt,draw=rulegray,text width=5.1cm,align=center,minimum height=13mm,inner sep=2mm},
 arrow/.style={->,thick,accent}
]
\node[start] (begin) at (0,0) {Сравниваются две точки или два вращающихся объекта};
\node[choice] (same) at (0,-2.15) {Точки принадлежат одному твёрдому телу?};
\node[action] (disk) at (-4.5,-5.10) {Одинаковы $\omega$, $\nu$, $T$\\$\displaystyle\frac{v_1}{v_2}=\frac{R_1}{R_2}$};
\node[choice] (gear) at (3.0,-5.05) {Шестерни сцеплены без проскальзывания?};
\node[note] (endleft) at (-4.5,-8.25) {Пользуйтесь отношением радиусов для линейных скоростей.};
\node[action] (contact) at (3.0,-8.10) {Равны модули скоростей на ободах\\$\displaystyle\frac{\nu_1}{\nu_2}=\frac{R_2}{R_1}$};
\node[note] (other) at (3.0,-11.15) {В остальных случаях используйте общие формулы $v=\omega R$, $\omega=2\pi\nu$.};
\draw[arrow] (begin)--(same);
\draw[arrow] (same.south west) -| node[above,pos=0.23] {да} (disk.north);
\draw[arrow] (same.south east) -| node[above,pos=0.23] {нет} (gear.north);
\draw[arrow] (gear)-- node[right] {да} (contact);
\draw[arrow] (gear.east) -- ++(0.75,0) |- node[right,pos=0.20] {нет} (other.east);
\draw[arrow] (disk)--(endleft);
\end{tikzpicture}
\end{center}
\vspace{5mm}
\begin{center}\small Один общий вопрос определяет, \emph{какое именно равенство} можно использовать.\end{center}

\clearpage
\topic{Тренировочный блок \texorpdfstring{\quad}{ }\textnormal{(Путь А)}}
Решите задания по порядку. В каждом случае запишите исходную формулу и необходимые преобразования. Используйте $\pi\approx3{,}14$, если требуется численная оценка.
\begin{enumerate}
 \item Точка прошла $\frac34$ окружности за $3\unit{с}$. Найдите частоту и период её равномерного движения.
 \item Колесо совершило $48$ оборотов за $12\unit{с}$. Найдите частоту и период.
 \item Период вращения равен $2{,}5\unit{с}$. Сколько оборотов совершит тело за $30\unit{с}$?
 \item Точка движется по окружности радиуса $0{,}5\unit{м}$ с периодом $4\unit{с}$. Найдите линейную и угловую скорости (точно через $\pi$).
 \item Частота вращения диска равна $5\unit{Гц}$. Найдите угловую и линейную скорости точки на расстоянии $20\unit{см}$ от оси.
 \item Полагая радиус Земли равным $6400\unit{км}$, а время одного оборота вокруг оси равным $24\unit{ч}$, оцените линейную скорость точки экватора в м/с. Землю считайте шаром.
 \item На одном жёстком диске точки $A$ и $B$ находятся на расстояниях $12\unit{см}$ и $30\unit{см}$ от оси. Если $v_A=0{,}8\unit{м/с}$, найдите $v_B$.
 \item Две шестерни радиусами $18\unit{см}$ и $12\unit{см}$ сцеплены без проскальзывания. Первая совершает $36$ оборотов за $9\unit{с}$. Найдите частоту вращения второй и укажите направление её вращения относительно первой.
 \item Тело движется равномерно по окружности радиуса $8\unit{м}$ со скоростью $12\unit{м/с}$. Найдите центростремительное ускорение. Во сколько раз оно изменится, если скорость увеличить втрое, а радиус --- вдвое?
 \item Две внешние шестерни радиусами $24\unit{см}$ и $12\unit{см}$ сцеплены без проскальзывания. Частота первой $3\unit{Гц}$. Найдите частоту второй, линейную скорость точки второй шестерни, расположенной в $6\unit{см}$ от её центра, и центростремительное ускорение этой точки (точно через $\pi$).
\end{enumerate}

\clearpage
\topic{Ответы к тренировочному блоку}
\renewcommand{\arraystretch}{1.42}
\begin{tabularx}{\linewidth}{@{}>{\bfseries\color{darkaccent}}p{1.15cm}X@{}}
\toprule
№ & Ответ \\
\midrule
1 & $\nu=0{,}25\unit{Гц}$; $T=4\unit{с}$. \\
2 & $\nu=4\unit{Гц}$; $T=0{,}25\unit{с}$. \\
3 & $n=12$ оборотов. \\
4 & $v=\dfrac{\pi}{4}\unit{м/с}$; $\omega=\dfrac{\pi}{2}\unit{рад/с}$. \\
5 & $\omega=10\pi\unit{рад/с}$; $v=2\pi\unit{м/с}$. \\
6 & $v\approx465\unit{м/с}$. \\
7 & $v_B=2\unit{м/с}$. \\
8 & $\nu_2=6\unit{Гц}$; вращение в противоположную сторону. \\
9 & $a_{\mathrm c}=18\unit{м/с}^2$; увеличится в $\dfrac92=4{,}5$ раза. \\
10 & $\nu_2=6\unit{Гц}$; $v=0{,}72\pi\unit{м/с}$; $a_{\mathrm c}=8{,}64\pi^2\unit{м/с}^2$. \\
\bottomrule
\end{tabularx}
\vspace{8mm}
\textbf{Самопроверка:} сумейте словами объяснить, почему у двух точек одного диска равны угловые скорости, а у двух сцепленных шестерён равны модули линейных скоростей на ободах.

\end{document}
